% GDZ PPN236018647,native544lower–545;HKUST545–546.
\NumberedItem{61}
\indent A potestatibus Integ[ralis]
\[
\int\frac{dx}{\surd(1-x^4)}\qquad(0\ldotp\,\ldotp\,\ldotp1)
\]
\par\vfill\noindent\hfill65\textsuperscript{*}

\clearpage\SourceLeaf{545}\ArithmeticHead{516}{L}{TAGEBUCH MIT ERLÄUTERUNGEN.}
\noindent pendet
\[
\Sigma\left(\frac{mm+6mn+nn}{(mm+nn)^4}\right)^k.
\]
\noindent\hfill[1797 Mart.]

\medskip\begingroup\small
\indent Es ist für positive ganzzahlige Werte von \(k\)
\[
S_k=\sum_{m,n}\left\{\frac{m^4-6m^2n^2+n^4}{(m^2+n^2)^4}\right\}^k
=\frac12\sum_{m,n}\left\{\frac1{(m+n\sqrt{-1})^4}+\frac1{(m-n\sqrt{-1})^4}\right\}^k,
\]
\noindent wo \(m,n\) alle Paare ganzer Zahlen mit Ausnahme von \(m=0,n=0\) durchlaufen. Die \(S_k\) lassen sich also in einfacher Weise durch die
\[
s_k=\sum_{m,n}\frac1{(m+n\sqrt{-1})^{4k}}\qquad(k=1,2,3,\ldotp\,\ldotp\,\ldotp)
\]
\noindent ausdrücken, und von dieser Summe läßt sich leicht zeigen, daß sie sich von der \(k\)-ten Potenz der Größe
\[
\Pi=\int_0^1\frac{dx}{\sqrt{1-x^4}}
\]
\noindent nur durch einen rationalen Zahlenfaktor unterscheidet. Man braucht dazu nur auf die Formeln zurückzugreifen, die \textsc{Gauss} um dieselbe Zeit (März 1797) in den \textsc{Leiste} eingetragen hat, und die in dem Abschnitt [II.] über \textit{Lemniskatische Funktionen}, oben S. 153 ff. abgedruckt sind. Man nehme nämlich von den einzelnen Faktoren des doppelt unendlichen Produkts, das in der Darstellung von \(\operatorname{sinlemn}x\), oben art. [2.], S. 153 den Zähler bildet, den Logarithmus, also \(\displaystyle\log\left(1-\frac1{(m+n\sqrt{-1})^4}\frac{x^4}{\Pi^4}\right)\), entwickle diesen nach Potenzen von \(\displaystyle\left\{\frac1{m+n\sqrt{-1}}\frac{x}{\Pi}\right\}^4\) und summiere in bezug auf \(m,n\) wie angegeben. Man erhält dann für den Logarithmus des Zählers von \(\operatorname{sinlemn}x\), abgesehen von dem Gliede \(\log x\), eine nach Potenzen von \(\dfrac{x^4}{\Pi^4}\) fortschreitende Reihe, in der der Koeffizient von \(-\dfrac1k\dfrac{x^{4k}}{\Pi^{4k}}\) gleich \(s_k\) ist. Vergleicht man diese Reihe mit der auf anderm Wege erhaltenen Entwickelung dieses Zählerlogarithmus (siehe oben art. [10.], S. 159 den »log. Numeratoris«), wo die Koeffizienten rationale Zahlen sind, so ergibt sich die Richtigkeit unserer Aussage. Man wird mit Sicherheit behaupten können, daß \textsc{Gauss} diesen Gedankengang verfolgt hat, da es sich so am ungezwungensten erklärt, weshalb \textsc{Gauss} damals gerade auch die Entwickelungen der Logarithmen von Zähler und Nenner des \(\operatorname{sinlemn}x\) aufgezeichnet hat; vergl. auch die sehr viel spätere Aufzeichnung (auf S. 88 des Mai 1809 begonnenen Handbuchs 19, Be), die Werke III, S. 408 abgedruckt ist. Allerdings würde dann das \(mm+6mn+nn\) in unserer Tagebuchnotiz als Schreibfehler für \(m^4-6m^2n^2+n^4\) gelten müssen, was aber umso eher angeht, als die Größe
\[
\frac{m^4-6m^2n^2+n^4}{(m^2+n^2)^4}
\]
\noindent in dem Doppelprodukt oben S. 153 wirklich vorkommt und auch die Summe \(S_1\) noch an zwei andern Stellen dieser Leisteaufzeichnung (oben S. 155, 156) erscheint. \textsc{Gauss} würde dann hier zum ersten Male in der Entwicklungsgeschichte der elliptischen Funktionen die Klasse von Reihen betrachtet haben, die in allgemeinster Form in \textsc{Eisensteins} Abhandlung \textit{Genaue Untersuchung der unendlichen Doppelproducte, aus welchen die elliptischen Functionen als Quotienten zusammengesetzt sind}\footnote{\textsc{Crelles} Journal für Mathematik \textbf{35}, 1847, S. 153, Mathematische Abhandlungen, 1847, S. 213.} auftreten.
\par\noindent\hfill\textsc{Klein.\quad Schlesinger.}
\par\endgroup
